\subsection{Elastic collision: impulse and torque}

For the elastic collision, the friction is neglected, so the panel can act on the fragment only along the normal, 
delivering an impulse $J\hat{\vect n}$. By Newton's
third law the satellite gets $-J\hat{\vect n}$. The impulse–momentum theorem, applied to each body, and its rotational version 
(torque $=$ lever arm $\times$ force), applied to the satellite. We get,
\begin{equation}
  m\,\Delta\vect v_d = +J\hat{\vect n},\qquad
  M\,\Delta\vect V = -J\hat{\vect n},\qquad
  I\,\Delta\omega = \left(\vect r\times(-J\hat{\vect n})\right)_z = -J\,b,
  \label{eq:impulse}
\end{equation}
\begin{itemize}
  \item $m$ and $M$ (in $\mathrm{kg}$) are the masses of the fragment and the
satellite
\item $\Delta\vect v_d$ and $\Delta\vect V$ (in $\mathrm{m\,s^{-1}}$) are
their velocity changes
\item $J$ (in $\mathrm{N\,s}$) is the impulse magnitude
\item $I$ (in$\mathrm{kg\,m^{2}}$) is the satellite's moment of inertia about its centre of
mass (CoM)
\item $\Delta\omega$ (in $\mathrm{rad\,s^{-1}}$) is its change in angular
velocity
\item $\vect r$ (in $\mathrm{m}$) is the position of the impact point
relative to the CoM
\item $b=(\vect r\times\hat{\vect n})_z$ (in $\mathrm{m}$) is
the lever arm of the impulse about the CoM. 
\end{itemize}
The point of the panel that is hit moves with
the velocity $\vect V+\vect\omega\times\vect r$, so the satellite pushes back 
in two ways: by the motion of a whole satellite and by its motion. Hence every unit of
impulse changes the relative normal speed between the fragment and the struck
point by
\begin{equation}
  k = \frac{1}{m}+\frac{1}{M}+\frac{b^{2}}{I},
  \label{eq:k}
\end{equation}
\begin{itemize}
  \item $k$ is in $\mathrm{kg^{-1}}$
  \item $m$ and $M$ (in $\mathrm{kg}$) are the masses of the fragment and the
satellite
  \item $b^2/I$ is $\mathrm{m^2/(kg\,m^2)}$.
\end{itemize}
 The three terms come from the
fragment, the satellite moving and the satellite rotating ($b$ appear in the torque and in $v=\omega b$ (in $\mathrm{m\,s^{-1}}$)). In an elastic collision the
relative normal speed is reserved, as in the one-dimensional case,
which gives

\begin{equation}
  J=\frac{2v_n}{k},
  \label{eq:J-elastic}
\end{equation}

\begin{itemize}
  \item $v_n$ (in $\mathrm{m\,s^{-1}}$) is the normal component of the approach
speed
\end{itemize}
When the impulse passes through the CoM
($b=0$), $1/k=mM/(m+M)$ and
Eqs.~\eqref{eq:impulse}--\eqref{eq:J-elastic} reduce to the textbook 1-D result
for a target initially at rest,
$v_1'=\dfrac{m_1-m_2}{m_1+m_2}\,v_1$ and
$v_2'=\dfrac{2m_1}{m_1+m_2}\,v_1$ \citep{ling2016}

\begin{itemize}
  \item $m_1,m_2$ (in $\mathrm{kg}$) are the two masses
  \item $v_1$ (in $\mathrm{m\,s^{-1}}$) is the initial
speed of body~1
  \item $v_1',v_2'$ (in $\mathrm{m\,s^{-1}}$) are the final speeds.
\end{itemize}
Because the fragment's sideways velocity is unaffected, the
fragment rebounds with most of its energy.